\documentclass[11pt]{article}
\usepackage[a4paper,margin=28mm]{geometry}
\usepackage[T1]{fontenc}
\usepackage{lmodern}
\usepackage{amsmath,amssymb,amsthm}
\usepackage{microtype}
\usepackage[colorlinks=true,linkcolor=blue,citecolor=blue,urlcolor=blue]{hyperref}
\newtheorem{theorem}{Theorem}
\newtheorem{lemma}[theorem]{Lemma}
\newtheorem{proposition}[theorem]{Proposition}
\newcommand{\R}{\mathbb R}
\newcommand{\E}{\mathbb E}
\newcommand{\rate}{\operatorname{rate}}
\newcommand{\rank}{\operatorname{rank}}
\newcommand{\tr}{\operatorname{tr}}
\newcommand{\dist}{\operatorname{dist}}
\newcommand{\med}{\operatorname{med}}
\newcommand{\ed}{e^{\mathrm{det}}}
\newcommand{\er}{e^{\mathrm{ran}}}
\title{The maximal gain of nonadaptive randomization}
\author{Kaiwen Jin\\{\small\href{mailto:yui.ui.siki@gmail.com}{yui.ui.siki@gmail.com}}}
\date{}
\begin{document}
\maketitle

\begin{abstract}
We determine the supremum of the improvement in polynomial convergence rate
due to randomization for approximation of bounded linear operators between
real Banach spaces using nonadaptive continuous linear measurements.
The supremum is one. We construct fixed infinite-dimensional operators whose
gains tend to one, with worst-case expected-norm error and at most \(n\)
measurements for every random seed. The construction combines a
finite-dimensional deterministic lower bound with a known nonadaptive
CountSketch estimate and a weighted \(\ell_1\) sum of blocks.
\end{abstract}

\section{Introduction}

Suppose that a bounded linear operator \(S:X\to Y\), between real Banach
spaces, must be approximated on the closed unit ball \(B_X\) from \(n\)
continuous linear measurements. A nonadaptive deterministic algorithm has
the form
\[
 A(x)=\phi(L_1x,\ldots,L_nx),\qquad L_j\in X',
\]
where the functionals are fixed and the reconstruction \(\phi:\R^n\to Y\)
is unrestricted. Its minimal worst-case norm error is denoted by
\(\ed_n(S)\). Using fewer than \(n\) measurements is allowed.

A nonadaptive randomized algorithm chooses the whole list of functionals
and its reconstruction independently of the input. For every seed \(\omega\)
it uses at most \(n\) measurements. We require the real-valued error
\((\omega,x)\mapsto\|Sx-A^\omega(x)\|_Y\) to be jointly measurable on the
seed space times the Borel unit ball. The minimal error is
\[
 \er_n(S)=\inf_A\sup_{x\in B_X}\E\|Sx-A^\omega(x)\|_Y.
\]
Thus the input supremum is outside the expectation; the error criterion is
the expected norm. All errors in this note use nonadaptive information.

For a bounded nonnegative sequence \(u=(u_n)_{n\ge1}\), set
\[
 \rate(u)=\sup\{\alpha\ge0:u_n=O(n^{-\alpha})\}.
\]
Write \(a(S)=\rate(\ed_n(S))\), \(b(S)=\rate(\er_n(S))\), and
\begin{equation}\label{eq:G}
 G=\sup_{\substack{X,Y\ {\rm real\ Banach},\ S\in\mathcal L(X,Y)\\
                   a(S)<\infty,\ b(S)<\infty}}
       \bigl(b(S)-a(S)\bigr).
\end{equation}
The operator and spaces in each rate are fixed as \(n\to\infty\).

Krieg, Novak and Ullrich proved the general upper bound \(G\le1\)
\cite[Corollaries 14--15]{KNU}. If either space is Hilbert, their bound
improves to \(1/2\). The general nonadaptive question remained between
\(1/2\) and \(1\); this interval is stated explicitly in
\cite[Section 10.2, equation (10.10)]{KU}. It is also the question posed in
AIM Problem 555 \cite{AIM555}. Results that permit adaptive randomized
measurements do not decide this comparison.

\begin{theorem}\label{thm:main}
The supremum in \eqref{eq:G} is \(G=1\).
\end{theorem}

The finite-dimensional estimates needed for the lower construction have
different dependence on dimension. A deterministic method has constant
error on \(\ell_p^m\to\ell_\infty^m\) when \(m\) is sufficiently large
relative to its rank. In contrast, a nonadaptive randomized sketch needs
only logarithmic dependence on \(m\) for a prescribed accuracy
\cite[Proposition 4.1]{KW}. A finite-dimensional example alone has infinite
asymptotic rates, so it cannot be used in \eqref{eq:G}. We place such
examples in a fixed infinite-dimensional operator. The \(\ell_1\) norm
across blocks controls the expected error for inputs occupying many blocks;
reading an initial group of blocks exactly makes the remaining sketch
budget affordable. The resulting gains approach one. We do not assert
that the supremum is attained by one operator.

\section{Finite-dimensional estimates}

We first give a lower bound that allows arbitrary deterministic
reconstruction. For a vector \(z\in\R^m\), \(\|z\|_r\) denotes its usual
\(\ell_r\) norm.

\begin{lemma}\label{lem:det}
Let \(q\ge4\) be an even integer, \(p=q/(q-1)\), and \(0\le n<m\).
Every deterministic algorithm using at most \(n\) fixed linear measurements
to approximate the identity \(\ell_p^m\to\ell_\infty^m\) has error at least
\begin{equation}\label{eq:finite-det}
 \frac{1-\sqrt{n/m}}{1+\sqrt n\,m^{-1/q}}.
\end{equation}
In particular its error is at least \(1/4\) if
\(m\ge4n\) and \(m\ge n^{q/2}\).
\end{lemma}

\begin{proof}
Let \(N\) be the information matrix and \(V\) its row space. Set
\[
 \rho=\max_{1\le i\le m}
       \sup_{\substack{x\in\ker N\\\|x\|_p\le1}}|x_i|.
\]
The inputs \(x\) and \(-x\) in \(\ker N\) have identical information.
Their images have distance \(2\|x\|_\infty\), so the worst error of any
reconstruction is at least \(\rho\).

For each standard basis vector \(e_i\), we claim that
\begin{equation}\label{eq:distance}
 \sup_{\substack{x\in V^\perp\\\|x\|_p\le1}}|x_i|
       =\dist_{\ell_q^m}(e_i,V).
\end{equation}
H\"older's inequality proves the inequality from left to right.
Choose a minimizer \(v_i\in V\) for the distance. Such a minimizer exists
by finite-dimensional compactness: a minimizing sequence is bounded and
\(V\) is closed. Put \(r_i=e_i-v_i\). If \(r_i=0\), both sides vanish.
Otherwise, differentiating the polynomial
\(\sum_j((r_i)_j-tv_j)^q\) at its minimum \(t=0\), for \(v\in V\), gives
\(\sum_j(r_i)_j^{q-1}v_j=0\). Therefore
\[
 x_j=\frac{(r_i)_j^{q-1}}{\|r_i\|_q^{q-1}}
\]
defines a vector in \(V^\perp\) with \(\|x\|_p=1\) and
\[
 x_i=\langle e_i,x\rangle
     =\langle r_i,x\rangle=\|r_i\|_q.
\]
This proves \eqref{eq:distance}.

Form the matrix \(B\) with rows \(v_i\), and put \(E=I-B\). Its residual
rows have \(\ell_q\) norm at most \(\rho\), and \(\rank B\le n\). Hence
\[
 \tr B\ge m(1-\rho),\qquad
 \|B\|_F\le\sqrt m+m^{1-1/q}\rho,
\]
where \(\|\cdot\|_F\) is the Frobenius norm. Indeed,
\(\|z\|_2\le m^{1/2-1/q}\|z\|_q\) for every row of \(E\).
Let \(P_V\) be the Euclidean orthogonal projection onto \(V\).
Since \(B=BP_V\), \(P_V\) is symmetric and
\(\|P_V\|_F=\sqrt{\dim V}\), we have
\[
 |\tr B|=|\langle B,P_V\rangle_F|\le\sqrt n\,\|B\|_F.
\]
Combining these inequalities and dividing by \(m\) yields
\[
 1-\rho\le\sqrt{n/m}+\sqrt n\,m^{-1/q}\rho,
\]
which is \eqref{eq:finite-det}. Under the two additional size conditions
its numerator is at least \(1/2\) and its denominator at most two.
\end{proof}

We use the following explicit randomized estimate from
\cite[Proposition 4.1]{KW}, with its fixed preprint version and numbering.
For integers \(R,G\), with \(R\) odd, draw independent uniform labels
\(H_{ri}\in\{1,\ldots,G\}\) and signs \(\sigma_{ri}\in\{-1,1\}\).
Take the \(RG\) grouped linear measurements
\[
 y_{rg}=\sum_{\{i:H_{ri}=g\}}\sigma_{ri}z_i,
 \qquad
 C_{R,G}(z)_i=\med_{1\le r\le R}
                   \bigl(\sigma_{ri}y_{r,H_{ri}}\bigr).
\]
The labels and signs determine every measurement before the input.

\begin{proposition}[Kunsch--Wnuk]\label{prop:sketch}
Let \(1\le p\le2\), \(m\ge2\), \(L\ge1\) be an integer and \(G=2^{4+L}\).
Choose \(R\) as the smallest odd integer satisfying
\(R\ge\max\{5,2+3\log_2m\}\). Then
\begin{equation}\label{eq:sketch}
 \sup_{\|z\|_p\le1}\E\|z-C_{R,G}(z)\|_\infty
       \le4\,2^{-L/p}.
\end{equation}
The algorithm uses \(RG\) measurements for every seed.
\end{proposition}

An odd median commutes with multiplication by any real scalar.
Thus \eqref{eq:sketch} holds for arbitrary \(z\) after multiplying its
right side by \(\|z\|_p\). The seed space is finite, and every realized
reconstruction is continuous. These facts will also give the required
measurability after the blocks are combined.

\section{Infinite-dimensional construction}

Fix an even integer \(q\ge4\), and put
\[
 p=\frac q{q-1},\qquad s=\frac1{q^2},\qquad
 m_k=2^k,\qquad w_k=2^{-sk}\quad(k\ge1).
\]
Define the real Banach spaces
\[
 X_q=\left(\bigoplus_{k\ge1}\ell_p^{m_k}\right)_{\ell_1},
 \qquad \|x\|_{X_q}=\sum_{k\ge1}\|x_k\|_p,
 \qquad
 Y=\ell_\infty\bigl(\{(k,i):k\ge1,\ 1\le i\le m_k\}\bigr),
\]
and the real-linear operator
\begin{equation}\label{eq:operator}
 (S_qx)_{k,i}=w_k(x_k)_i.
\end{equation}
The direct sum is complete: a Cauchy sequence converges in each block;
passing to the limit in finite partial sums shows that its coordinate
limit belongs to \(X_q\) and that convergence holds in the summed norm.
The same coordinate-limit argument proves completeness of \(Y\).
Moreover \(\|S_qx\|_Y\le\|x\|_{X_q}\).

\begin{theorem}\label{thm:rates}
For the fixed operator \eqref{eq:operator}, both rates are finite and
\begin{equation}\label{eq:rates}
 s\le a(S_q)\le\frac{sq}{2},
 \qquad
 s+\frac1p\le b(S_q)\le a(S_q)+1.
\end{equation}
In particular,
\begin{equation}\label{eq:gain}
 b(S_q)-a(S_q)\ge1-\frac3{2q}+\frac1{q^2}.
\end{equation}
\end{theorem}

\begin{proof}
For the deterministic lower bound, let \(n\ge4\) and choose
\(k=\lceil(q/2)\log_2n\rceil\). Restrict an arbitrary \(n\)-measurement
algorithm on \(X_q\) to inputs in block \(k\), and project its output to
that block. Its restricted information has rank at most \(n\).
Since \(m_k\ge n^{q/2}\ge4n\), Lemma~\ref{lem:det} gives
\begin{equation}\label{eq:global-det}
 \ed_n(S_q)\ge\frac{w_k}{4}
       \ge\frac{2^{-s}}4\,n^{-sq/2}.
\end{equation}
Thus \(a(S_q)\le sq/2\). Conversely, reading all coordinates in blocks
\(1,\ldots,J\) costs \(2^{J+1}-2\) measurements and leaves error
\(w_{J+1}\). Taking \(J=\lfloor\log_2(n+2)\rfloor-1\), for \(n\ge2\),
gives \(\ed_n(S_q)\le2^s(n+2)^{-s}\), hence \(a(S_q)\ge s\).

We now allocate a strict randomized budget. In a block of dimension
\(m_k=2^k\), suppose \(w_k\ge\varepsilon>0\). Apply
Proposition~\ref{prop:sketch} with
\[
 L_k=\left\lceil p\log_2\frac{4w_k}{\varepsilon}\right\rceil .
\]
Then \(L_k\ge1\), the weighted expected error is at most
\(\varepsilon\|x_k\|_p\), and its number of measurements satisfies
\begin{equation}\label{eq:block-cost}
 R_kG_k\le 6k\cdot32
       \left(\frac{4w_k}{\varepsilon}\right)^p
       \le C_0k\left(\frac{w_k}{\varepsilon}\right)^p,
 \qquad C_0=3072.
\end{equation}
Here \(R_k\le3+3k\le6k\) and \(1<p\le2\).
All group measurements extend to continuous linear functionals on \(X_q\),
since each has support in one finite block.

Put \(\theta=2^{-sp}\), so \(0<\theta<1\), and fix \(D>0\) such that
\begin{equation}\label{eq:D}
 D^p\ge\frac{2C_0\,8^{sp}}{(\ln2)(1-\theta)^2}.
\end{equation}
For \(n\ge8\), set
\[
 J=\lfloor\log_2(n/4)\rfloor,\qquad
 \varepsilon=D\,n^{-s-1/p}\bigl(\ln(n+1)\bigr)^{1/p},\qquad
 K=\max\left\{J,\left\lfloor\frac{\log_2(1/\varepsilon)}s
                                      \right\rfloor\right\}.
\]
Read blocks \(k\le J\) exactly, sketch blocks \(J<k\le K\), and return
zero on the remaining blocks. These choices depend only on \(n,q\).
Every sketched block satisfies \(w_k\ge\varepsilon\), while
\(w_{K+1}<\varepsilon\); if \(K=J\), there are no sketched blocks.

The exact prefix costs at most \(n/2\). For the other measurements,
the geometric first-moment sum gives
\[
 \sum_{k>J}k\theta^k
 =\theta^{J+1}\left(\frac{J+1}{1-\theta}
                           +\frac{\theta}{(1-\theta)^2}\right)
 \le\frac{(J+1)\theta^J}{(1-\theta)^2}.
\]
Using \(2^J\ge n/8\) and \(J+1\le\log_2n\), \eqref{eq:block-cost} yields
\[
 \begin{split}
 \sum_{k=J+1}^{K}R_kG_k
 &\le C_0\varepsilon^{-p}\sum_{k>J}k2^{-spk}\\
 &\le\frac{C_0\,8^{sp}}{D^p(\ln2)(1-\theta)^2}
                    \frac{\ln n}{\ln(n+1)}\,n
 \le\frac n2.
 \end{split}
\]
Thus every seed uses at most \(n\) measurements. There are finitely many
sketched blocks, with a finite product seed space, and all their grouped
measurements and reconstructions are fixed before input.

For any \(x\in X_q\), bound the maximum block error by the sum of
the finitely many sketched errors and the deterministic tail bound.
Homogeneity of the block algorithms gives
\[
 \begin{split}
 \E\|S_qx-A_n^\omega(x)\|_Y
 &\le\varepsilon\sum_{J<k\le K}\|x_k\|_p
       +\varepsilon\sum_{k>K}\|x_k\|_p\\
 &\le\varepsilon\|x\|_{X_q}.
 \end{split}
\]
This includes inputs spread over infinitely many blocks. Each realized
algorithm is continuous, because its output consists of finitely many
linear coordinates and medians. Its scalar norm error is continuous in
\(x\), and hence jointly measurable on the finite seed space times \(B_{X_q}\).
We have proved
\begin{equation}\label{eq:global-ran}
 \er_n(S_q)\le D\,n^{-s-1/p}\bigl(\ln(n+1)\bigr)^{1/p},
 \qquad n\ge8.
\end{equation}
For every \(\alpha<s+1/p\), the logarithmic factor is
\(O(n^{s+1/p-\alpha})\), so \(b(S_q)\ge s+1/p\).

It remains to check the finite-rate restriction. The symmetric case of
\cite[Corollary 14]{KNU} gives, for every bounded linear \(S:X\to Y\)
and \(\alpha>0\),
\begin{equation}\label{eq:universal}
 \ed_{2n-1}(S)\le C_\alpha n^{1-\alpha}
          \sup_{0\le j<n}(j+1)^\alpha e_j^{\mathrm{ran,ad}}(S),
\end{equation}
where \(e_j^{\mathrm{ran,ad}}\) allows adaptive randomized information.
Its continuous-information formulation is explicit in
\cite[Section 6.3]{KNU}. Since
\(e_j^{\mathrm{ran,ad}}(S)\le\er_j(S)\), an available nonadaptive
randomized exponent \(\alpha>1\) implies \(a(S)\ge\alpha-1\).
The \(j=0\) initial error is bounded by \(\|S\|\), and monotonicity in the
deterministic budget extends the estimate to all indices.
If \(b(S_q)=\infty\), choosing an available \(\alpha>1+sq/2\)
contradicts \eqref{eq:global-det}. Hence \(b(S_q)<\infty\). If
\(b(S_q)>1\), taking \(\alpha\) up to \(b(S_q)\) proves
\(b(S_q)\le a(S_q)+1\); if \(b(S_q)\le1\), this follows from
\(a(S_q)\ge0\).
Finally \eqref{eq:rates} gives
\[
 b(S_q)-a(S_q)\ge s+\frac1p-\frac{sq}{2}
                    =1-\frac3{2q}+\frac1{q^2}.
\]
\end{proof}

\begin{proof}[Proof of Theorem~\ref{thm:main}]
The comparison \eqref{eq:universal} gives \(b(S)-a(S)\le1\) whenever
both rates are finite. Indeed, if \(b(S)>1\), use each available
\(1<\alpha<b(S)\) and then let \(\alpha\uparrow b(S)\); if \(b(S)\le1\),
use \(a(S)\ge0\). Thus \(G\le1\).
The operators of Theorem~\ref{thm:rates} are admissible in \eqref{eq:G}
for every finite even \(q\ge4\), and their gains in \eqref{eq:gain}
tend to one as \(q\to\infty\). Therefore \(G\ge1\).
\end{proof}

\paragraph{AI-use disclosure.}
This manuscript was developed with assistance from OpenAI Codex in exploring
the construction, drafting and checking the argument, verifying cited sources,
and preparing the text.

\begin{thebibliography}{9}
\bibitem{KNU}
D. Krieg, E. Novak and M. Ullrich,
\emph{On the power of adaption and randomization},
Forum of Mathematics, Sigma \textbf{13} (2025), e152.
\href{https://doi.org/10.1017/fms.2025.10101}{doi:10.1017/fms.2025.10101}.
Theorem numbering here refers to
\href{https://arxiv.org/abs/2406.07108v2}{arXiv:2406.07108v2}.

\bibitem{KW}
R. J. Kunsch and M. Wnuk,
\emph{Adaptive and non-adaptive randomized approximation of
high-dimensional vectors},
Journal of Approximation Theory \textbf{318} (2026), 106318.
\href{https://doi.org/10.1016/j.jat.2026.106318}{doi:10.1016/j.jat.2026.106318}.
Theorem numbering here refers to
\href{https://arxiv.org/abs/2410.23067v2}{arXiv:2410.23067v2}.

\bibitem{KU}
D. Krieg and M. Ullrich,
\emph{Approximation of functions: Optimal sampling and complexity},
Acta Numerica \textbf{35} (2026), 273--457.
\href{https://doi.org/10.1017/S0962492925100287}{doi:10.1017/S0962492925100287}.
The equation locator here refers to
\href{https://arxiv.org/abs/2602.02066v2}{arXiv:2602.02066v2}.

\bibitem{AIM555}
AIM problem collection,
\emph{Problem 555: Maximal gain from randomization with nonadaptive
linear measurements},
statement at revision \texttt{c9929805cece044705abaaebe1cebddac1ba11a6}.
\href{https://github.com/MColbrook/AIM/blob/c9929805cece044705abaaebe1cebddac1ba11a6/problems/555-nonadaptive-randomization-gain.md}{Immutable problem statement}.
\end{thebibliography}
\end{document}
